% Shared preamble for all lab reports (Tufte handout layout).
%
% Usage, from labN/report/labN_report.tex:
%   \documentclass[a4paper,justified]{tufte-handout}
%   \input{../../report-style/preamble}
%   \labtitle{2}{ROS, TF and Quaternions}
%
% Build: latexmk -pdf labN_report.tex

\usepackage{amsmath,amssymb,bm}
\usepackage{booktabs,tabularx,array}
\usepackage{graphicx}
\usepackage{xcolor}
\usepackage{textcomp}
\usepackage{listings}
\usepackage[most]{tcolorbox}
\usepackage{siunitx}
\usepackage{cleveref}

\setcounter{secnumdepth}{2}
\crefname{section}{Section}{Sections}
\crefname{figure}{Figure}{Figures}
\crefname{table}{Table}{Tables}
\crefname{listing}{Listing}{Listings}
% In main-column table/figure floats a \label after \caption refers to the
% section (tufte-latex with the current kernel); write
% \caption{Text.\protect\label{key}} instead. Margin floats are unaffected.
\hypersetup{colorlinks=true, allcolors=[rgb]{0.12,0.25,0.55}}

% --- Author and course -------------------------------------------------------
\newcommand{\studentname}{Xuanye Xu}
\newcommand{\studentids}{24020036071 / H00490299}
\newcommand{\coursename}{Robotics Integration I --- VNAV Labs}

% \labtitle{<lab number>}{<topic>}
\newcommand{\labtitle}[2]{%
  \title[Lab #1: #2]{Lab #1 Report: #2}%
  \author[\studentname]{\studentname\\[2pt]
    {\normalsize\normalfont Student ID \studentids\\ \coursename}}%
  \date{\today}}

% --- Grading items -----------------------------------------------------------
% \gradeditem{<label>}{<title>}{<points>} starts a section for one grading item.
\newcommand{\gradeditem}[3]{%
  \section{#2\texorpdfstring{\enspace{\normalfont\small[#3~pts]}}{ [#3 pts]}}\label{#1}}

% Evidence summary shown at the end of each graded section.
\newtcolorbox{evidence}{enhanced, breakable, frame hidden, arc=0pt,
  colback=black!4, borderline west={1.2pt}{0pt}{black!55},
  left=6pt, right=4pt, top=3pt, bottom=3pt,
  before upper={\textsc{Evidence}\quad}}

% Marks unfinished content in red; a final report must contain none.
\newcommand{\pending}[1]{\textcolor{red}{[PENDING: #1]}}

% --- Code and terminal output ------------------------------------------------
\lstdefinestyle{code}{language=C++, basicstyle=\ttfamily\footnotesize,
  keywordstyle=\bfseries, commentstyle=\itshape\color{black!60},
  numbers=left, numberstyle=\tiny\color{black!50}, xleftmargin=1.5em,
  columns=fullflexible, keepspaces, breaklines, showstringspaces=false,
  upquote=true}
\lstdefinestyle{log}{basicstyle=\ttfamily\scriptsize, frame=leftline,
  framerule=0.6pt, rulecolor=\color{black!40}, xleftmargin=0.8em,
  columns=fullflexible, keepspaces, breaklines, upquote=true}
\lstset{style=log}
